%% =================================================================
%% Wei, Mei, Zhao, Xiao, Sun, Zhang, Guo.
%% "Nondestructively identifying the mechanical behavior of soft
%%  tissues using surface deformation with an explicit inverse
%%  approach."
%% Appl. Math. Modelling 134 (2024) 126--147.
%%
%% Reconstruction of page 18 (printed p. 143) as study notes.
%% Generated by: reproduce-basic-paper-tex → reproduce-multiple-pages-basic-tex [17, 22]
%%
%% Structure: 4 prose paragraphs (finishing Case 3 discussion) +
%% Discussion subsection + Fig. 24 (Sample 1 with 1% noise).
%% No equations, no tables.
%% =================================================================

\documentclass[11pt]{article}

\usepackage[margin=1in]{geometry}
\usepackage{amsmath, amssymb}
\usepackage{mathrsfs}
\usepackage{bm}
\usepackage{xcolor}
\usepackage{fancyhdr}

\pagestyle{fancy}
\fancyhf{}
\lhead{\small Y.~Mei et al.}
\rhead{\small Applied Mathematical Modelling 134 (2024) 126--147}
\cfoot{\small 143 $\to$ \arabic{page}}
\renewcommand{\headrulewidth}{0.4pt}

\setcounter{page}{1}

\begin{document}

\noindent
while the implicit inverse method required a significantly higher
number of iterations, reaching convergence after 4354 iterations.

In the absence of noise in the datasets, both the explicit and
implicit inverse methods exhibited excellent performance in
yielding shear modulus reconstructions, as demonstrated in
Fig.~23. However, when 1\,\% noise was introduced to the
datasets, the explicit inverse method outperformed the implicit
inverse method (see Fig.~24). For the explicit inverse method,
the average relative error in the reconstructed shear modulus
values was approximately 2\,\%, indicating a relatively accurate
estimation despite the presence of noise (see Table~5). On the
other hand, the implicit inverse method exhibited a higher
average error of more than 3\,\%, signifying a less precise
reconstruction in the presence of noise. Overall, the explicit
inverse method proved to be more robust and capable of handling
noisy datasets, leading to more accurate shear modulus
reconstructions compared to the implicit inverse method.
Furthermore, it is noteworthy that the reconstructed results
maintain good quality even when the noise level is increased to
5\,\% (see Fig.~25).

In the case where two inclusions were embedded in the background
(as depicted in Fig.~21(b)). As shown in Fig.~26, the results
obtained using the implicit inverse method showed a consistent
underestimation of the recovered shear modulus values by at
least 20\,\% as the size of the mapped inclusions increased.
Additionally, the shape of the mapped inclusions appeared to
become more circular with the increase of surface displacement
datasets. Conversely, when employing the explicit inverse
method, the shear modulus values within the two inclusions were
overestimated by approximately 20\,\%. Furthermore, the shape
of the recovered inclusions exhibited a more circular appearance
compared to those predicted by the implicit inverse method. The
regularization factors used in Figs.~23--26 are listed in
Table~6.

To assess the robustness of the proposed explicit inverse
method, we varied the shape of the inclusion in Sample~1, as
depicted in Fig.~27(a) and~(d). Our observations indicate that
the reconstructed results obtained from the explicit inverse
method effectively preserve the shape of the complex inclusion,
even with only three displacement datasets. In contrast, the
implicit inverse method failed to achieve this level of
accuracy.

\vspace{0.5em}

\noindent\textbf{Discussion}

\vspace{0.3em}

\noindent
The paper proposes an explicit inverse approach for
characterizing the nonhomogeneous elastic property distribution
of three typical nonlinear biological structures with surface
deformation. This approach is different from the prevalent
implicit inverse approach in that it optimizes only the
geometric parameters and shear moduli of each region, reducing
the scale of the optimization problem. By doing so, the proposed
method improves the uniqueness of the inverse solution since the
total number of optimization variables is reduced. Additionally,
the paper highlights an important advantage of the explicit
method: it only requires surface deformation data, which can be
measured with high accuracy using low-cost measurement
techniques. This makes the proposed approach more practical and
accessible for real-world applications, as it avoids the need
for complex and expensive instrumentation to gather the
full-field data.

\vspace{0.8em}

% ---------- Figure 24 placeholder ----------
\noindent
\begin{center}
\fbox{\parbox{0.92\textwidth}{\centering\vspace{1em}
\textbf{[Figure 24 --- placeholder]}\\[0.8em]
Six-panel comparison of \textbf{Explicit Inverse Method} (top
row) vs.\ \textbf{Implicit Inverse Method} (bottom row) on
\textbf{Case 3 Sample 1} (one embedded inclusion, target
$\approx 5$\,MPa), with \textbf{1\,\% noise} added, using 3, 6,
and 9 surface displacement datasets.\\[0.4em]
\textbf{Explicit (top row)}:
(a) 3 fields, SM $\in [1.00, 5.75]$\,MPa --- \textbf{15\,\%
overshoot} of the 5\,MPa nominal tumor stiffness;
(b) 6 fields, SM $\in [1.00, 5.30]$\,MPa;
(c) 9 fields, SM $\in [1.00, 5.01]$\,MPa --- near-perfect with
9 fields. The max drifts downward monotonically
($5.75 \to 5.30 \to 5.01$) toward the nominal target.\\[0.4em]
\textbf{Implicit (bottom row)}:
(d) 3 fields, SM $\in [1.00, 3.26]$\,MPa --- \textbf{35\,\%
undershoot}, a dramatically worse number than the paper's prose
claim of ``$\approx 3$\,\% average error'' suggests (the prose
error is measured over the whole domain, not the inclusion
max);
(e) 6 fields, SM $\in [1.00, 3.51]$\,MPa --- 30\,\% undershoot;
(f) 9 fields, SM $\in [1.00, 4.82]$\,MPa --- still 3.6\,\% below
nominal. Max drifts upward monotonically
($3.26 \to 3.51 \to 4.82$) toward the target but doesn't
reach it.\\[0.3em]
\textbf{Takeaway}: with 1\,\% noise, the explicit and implicit
methods fail in \emph{opposite directions} on this problem ---
explicit \emph{overshoots} the tumor stiffness, implicit
\emph{undershoots} it. Both approach the 5\,MPa target from
their respective sides as more fields are added, but the
explicit method gets closer faster. The prose's ``2\,\% vs
3\,\% average error'' framing understates the practical
significance: a 35\,\% undershoot of the tumor stiffness
(panel d) would have very different clinical implications than
a 2\,\% average-domain error suggests.
\vspace{1em}}}
\end{center}
\vspace{0.3em}
\begin{center}
\textbf{Fig. 24.} The reconstructed results with varying 3, 6
and 9 surface displacement datasets corresponding to Sample 1
(1\,\% noise is introduced into the measured data, Unit: MPa).
\end{center}

% [page ends here — continues on page 19]

\end{document}
